\textbf{1. Kiến thức cần nhớ}
\begin{itemize}
\item \textbf{Tần số tương đối} của một giá trị là tỉ số giữa tần số của giá trị đó và cỡ mẫu:
$$\boxed{f_i = \dfrac{m_i}{N} \times 100\%}$$
\item Tổng các tần số tương đối của tất cả các giá trị luôn bằng $100\%$:
$$\boxed{f_1 + f_2 + \dots + f_k = 100\%}$$
\item \textbf{Biểu đồ tần số tương đối}: Thường dùng biểu đồ cột hoặc biểu đồ hình quạt tròn (tần số tương đối $f_i\%$ ứng với góc ở tâm $\alpha_i = 360^\circ \times \dfrac{f_i}{100}$).
\end{itemize}

\textbf{2. Phương pháp giải các dạng}

%% DANG 9D7?23-1
\textbf{Dạng 1. Lập bảng tần số tương đối}
\begin{enumerate}
\item Tính cỡ mẫu $N = \sum m_i$.
\item Lấy từng tần số $m_i$ chia cho $N$ rồi nhân với $100\%$ để được $f_i\%$.
\item Kiểm tra tổng các $f_i$ có đúng bằng $100\%$ không (chú ý làm tròn số hợp lí).
\end{enumerate}

%% DANG 9D7?23-2
\textbf{Dạng 2. Vẽ và đọc biểu đồ tần số tương đối (cột, hình quạt tròn)}
\begin{enumerate}
\item Vẽ biểu đồ hình quạt tròn:
\begin{itemize}
\item Tính số đo góc ở tâm tương ứng: $\alpha_i = 360^\circ \cdot \dfrac{m_i}{N} = 3{,}6^\circ \times f_i$.
\item Dùng thước đo góc vẽ các hình quạt và ghi tỉ lệ phần trăm kèm chú giải.
\end{itemize}
\item Đọc biểu đồ: So sánh cơ cấu phần trăm, xác định nhóm chiếm tỉ lệ áp đảo.
\end{enumerate}
